\section{The language ILD}

\subsection{Syntax}
\label{syntax}
\begin{cfig}{Grammar}{Syntax of ILD}
\begin{tabular}{r@{\;}c@{\;}l@{\qquad}l}
\(\ildsetVal\) & \(:=\) & \(\ildsetSym \mid \ildsetSexp \mid \ildsetSF \mid \ildsetFail \mid \ildsetHost\) & \paraphrase{Tangible} values \\
\(\ildsetSym\) & \(\in\) & \emph{symbols} & \\
\(\ildsetSexp\) & \(:=\) & \(() \mid \ildpair{\ildsetVal}{\ildsetVal}\) & S-expressions \\
\(\ildsetSF\) & \(:=\) & \(\ildsf{free-vars} \mid \ildsf{quote} \mid \ildsf{macroexpand}\) & Special forms \\
\(\ildsetFail\) & \(:=\) & \(\ildfail{\ildsetVal}\) & Failure objects, carrying a context value \\
\(\ildsetHost\) & \(:=\) & \(\ildsetNum \mid \ildsetBool \mid \ildsetStr \mid \ildsetHostFunc \mid \dots\) & Host values and \\
\(\ildsetHostFunc\) & \(:=\) & \(\ildsetAbstr \mid \ildsetVCont \mid \ildhost{+} \mid \ildhost{cons} \mid\) & host functions, opaque \\
 & & \(\ildhost{apply} \mid \ildhost{call/cc} \mid \dots\) & to the base language (\cref{embedding}) \\
\(\ildsetNum\) & \(\in\) & \emph{numbers} & \\
\(\ildsetBool\) & \(:=\) & \(\ildhost{t} \mid \ildhost{f}\) & \\
\(\ildsetStr\) & \(\in\) & \emph{strings} & \\
\(\ildsetProg\) & \(:=\) & \(\ildsetSym \mid () \mid \ildpair{\ildsetProg}{\ildsetProg} \mid\) & Surface syntax (programs) \\
 & & \(\ildsetStr \mid \ildsetNum \mid \ildsetBool\) & \\
\(\ildsetList\) & \(:=\) & \(() \mid \ildpair{\ildsetVal}{\ildsetList}\) & S-expression \emph{lists} \\
\(\ildsetEnv\) & \(:=\) & \(() \mid \ildpair{\ildpair{\ildsetSym}{\ildsetVal}}{\ildsetEnv}\) & Binding environments\footnote{We assume keys are restricted to be unique (\cref{stepped-semantics}). Furthermore, performant implementations will use other representations of binding environments.} (K/V lists) \\
\(\ildsetAbstr\) & \(:=\) & \(\ildabstr{\ildsetEnv, \ildsetParList, \ildsetVal}\) & Abstraction terms, \cref{abstraction} \\
\(\ildsetParList\) & \(:=\) & \(() \mid \ildpair{\ildsetSym}{\ildsetParList}\) & Param lists (lists of symbols) \\
\(\ildsetVCont\) & \(:=\) & \(\ildcont{\ildsetCont}\) & Reified continuations, \cref{first-class-continuations} \\
\(\ildsetComp\) & \(:=\) & \(\interop{eval}{\ildsetCont}{\ildsetEnv}{\ildsetCVal} \mid\) & CPS calculus computation terms \\
 & & \(\interop{apply}{\ildsetCont}{\ildsetHostFunc}{\ildsetCVal^*} \mid\) & \\
 & & \(\interop{comb}{\ildsetCont}{\ildsetEnv}{\ildsetCVal}{\ildsetCVal^*} \mid\) & \\
 & & \(\cpsapp{\ildsetCont, \ildsetCVal}\) & \\
\(\ildsetCont\) & \(:=\) & \(\cpsyield \mid \cpsabstr{\ildsetVV}{\ildsetComp}\) & CPS calculus continuations \\
\(\ildsetCVal\) & \(:=\) & \(\ildsetVV \mid \ildsetVal\) & CPS calculus value terms \\
\(\ildsetVV\) & \(\in\) & \(\{x_1, x_2, \dots\}\) & CPS calculus value variables \\
\end{tabular}
\end{cfig}

We define ILD as a homoiconic language where parseable programs (\(\ildsetProg\))
are a subset of the internal syntax \(\ildsetVal\). Source code is parsed as
standard S-expressions. Throughout this paper, we use the \(\ildlist{a_1\, a_2\, \dots\, a_n}\)
syntax to denote S-expression lists, i.e. the nested pairs
\(\ildpair{a_1}{\ildpair{a_2}{\,\dots\, \ildpair{a_n}{()} \,\dots\,}}\).

In addition, the parser of ILD is augmented with the following syntax sugars:
\begin{itemize}
  \item Quote: \ild{'<expr>} \(\Mapsto\) \ild{(quote <expr>)}
  \item Macroexpand: \\
    \ild{(!<expr1> ... <exprN>)} \(\Mapsto\) \\
    \ild{(macroexpand <expr1> ... <exprN>)}
\end{itemize}

We also introduce a set of computation terms (\(\ildsetComp\)) that facillitate
the operational semantics.

\subsection{Semantics}
\label{semantics}
We define the semantics of ILD in terms of a CPS calculus based on
standard \(\beta\)-reduction%
\footnote{
We omit details related to variable renaming and substitution.
}%
over the set of \emph{computations} (\(\ildsetComp\)). Evaluation starts from
the term \(\interop{eval}{\cpsyield}{E}{v}\) for some value \(v\), the root
continuation \(\cpsyield\) and a root binding environment \(E\) (\cref{root-env}).

\subsubsection{Embedding}
\label{embedding}
ILD is designed to be embedded in a host environment that provides a set of datastructures
and library functions (\(\ildsetHost\)). We split the semantics of ILD into:
\begin{itemize}
  \item Base language, whose semantics are shown in this section
  \item FFI semantics of the host environment (\(\ildsetHostFunc\)) via
    the term \(\interopword{apply}\). A small host environment, sufficient
    for writing nontrivial programs, is discussed in \cref{root-env}.
\end{itemize}

\subsubsection{Notes on small-step semantics}
\label{stepped-semantics}
\begin{cfig}{Figure}{Small-step semantics of ILD}
\label{stepped-semantics-fig}
\begin{semgroup}{Simple cases:}
\semrow{%
  \step{s \in \ildsetSym \text{ and } e = \ildlist{\dots\, \ildpair{s}{v}\, \dots}}{\interop{eval}{C}{e}{s}}{\cpsapp{C,\, v}}
  \qquad
  \step{v \in \ildsetSF \cup \ildsetHost \cup \ildsetFail}{\interop{eval}{C}{e}{v}}{\cpsapp{C,\, v}}
}%
\end{semgroup}

\begin{semgroup}{Combinations:}
\semrow{%
  \evalsto{\interop{eval}{C}{e}{\ildlist{\phi\, \alpha_1\, \dots\, \alpha_n}}}{\interop{eval}{\cpsabstr{f}{\interop{comb}{C}{e}{f}{\alpha_1}{\dots}{\alpha_n}}}{e}{\phi}}
}%
\semrow{%
  \step{f \in \ildsetHostFunc}{\interop{comb}{C}{e}{f}{\alpha_1}{\dots}{\alpha_n}}{\interop{eval}{\cpsabstr{x_1}{\cdots \interop{eval}{\cpsabstr{x_n}{\interop{apply}{C}{f}{x_1}{\dots}{x_n}}}{e}{\alpha_n} \cdots}}{e}{\alpha_1}}
}%
\end{semgroup}

\begin{semgroup}{Special forms:}
\semrow{%
  \evalsto{\interop{comb}{C}{e}{\ildsf{quote}}{\alpha}}{\cpsapp{C,\, \alpha}}
  \qquad
  \evalsto{\interop{comb}{C}{e}{\ildsf{free-vars}}}{\cpsapp{C,\, e}}
}%
\semrow{%
  \evalsto{\interop{comb}{C}{e}{\ildsf{macroexpand}}{\phi}{\vec{\alpha}}}{\interop{eval}{\cpsabstr{f}{\interop{apply}{\cpsabstr{x}{\interop{eval}{C}{e}{x}}}{f}{\vec{\alpha}}}}{e}{\phi}}
}%
\end{semgroup}

\begin{semgroup}{CPS calculus \(\beta\)-reduction:}
\semrow{\evalsto{\cpsapp{\cpsabstr{x}{M},\, v}}{M\subst{x}{v}}}%
\end{semgroup}
\end{cfig}

Small-step semantics is outlined in \cref{stepped-semantics-fig}. A few notes
on selected cases:
\begin{itemize}
  \item \(\ildsf{free-vars}\) captures the current binding environment and is used by higher-level constructs
    such as the lambda macro (\cref{lambda-macro}).
  \item The head of a combination is always evaluated. If the result of that is a special
    form, the special form is applied on the \textbf{unevaluated} remaining elements of the combination.
    If the head is a host function, the operands are then \textbf{evaluated} and then the head is applied
    to the resulting arguments.
  \item \(\ildlist{\ildsf{macroexpand}\, \phi\, \alpha_1\, \alpha_2\, \dots\, \alpha_n}\) evaluates just \(f\) and then passes the
    \textbf{unevaluated} arguments to it. The result is then in turn evaluated. This mechanism
    is discussed in \cref{macroexpand-mechanism}.
  \item Congruence rules are omitted for brevity. Any \(\ildsetCont\) or \(\ildsetComp\) subterm
    is considered to be a reduction hole.
  \item All unlisted cases result in a \(\ildfailbare\).
\end{itemize}
